\scanpage{261}
\mspage{73}
% archive sheet number: 261
\uline{Construction des diagrammes} \uline{substantiels} pour un dérivateur

Soit $\mathbb D$ un prédérivateur, $I$ dans $\mathrm{Dia}$, considérons

\[
\varphi_I:\quad\mathbb D(I)\longrightarrow\underline{\operatorname{Hom}}(I^{\unsure{\mathrm o|\circ}},\mathcal A)
\qquad\text{où }\mathcal A=\unsure{\mathbb D(e)|D(e)}
\]

est la cat. des coefficients fondamentale. Les objets de $\underline{\operatorname{Hom}}(I^{\mathrm o},\mathcal A)$ s'interprètent comme des \uline{diagrammes de type} $I^{\mathrm o}$ dans $\mathcal A$. \inserted{\unsure{Ce langage est surtout commode} si $I$ est \unsure{ordonnée|ordonné}.} Notons, pour $F\in\mathrm{Ob}\,\mathbb D(I)$, par \deleted{\illegible{1}} $(F_i)_{i\in\mathrm{Ob}\,I^{\deleted{\unsure{\mathrm o}}}}$ le $I^{\mathrm o}$-diagramme associé dans $\mathcal A$. Par l'axiome Dér 2, le foncteur $\varphi$

\[
F\longmapsto(F_i)_{i\in\mathrm{Ob}\,I}
\]

est conservatif. Mais il n'a aucune envie d'être fidèle, encore moins pleinement fidèle. Quand on connaît \deleted{\illegible{1}}\unsure{le} diagramme $(F_i)_i$, on est loin de connaître $F$, pas \unsure{$=$|même} à \unsure{isom} non unique près je présume. Il n'est nullement clair que si $F,\deleted{\unsure{E}}\unsure{F'|E'}$ de $\mathbb D(I)$ sont tels que les diagrammes associés soient isomorphes, ils soient eux-mêmes isomorphes.

Pour insister sur ce point, les objets de $\mathbb D(I)$ s'appelleront les ``\uline{diagrammes substantiels} de type $I^{\mathrm o}$ dans $\mathcal A$'' (\unsure{ou} $\mathbb D$-substantiels, s'il y a risque de confusion). Pour étayer l'intuition et simplifier les explications, un objet $F$ de $\mathbb D(I)$ (\inserted{$\mathbb D$-}coefficient sur $I$) sera représenté graphiquement par le diagramme \inserted{``faible'' ou ``\unsure{squelettique}''} correspondant --- sous les regards
